\bookchapter{الفصل 56: الصقل والنشر — الاختبار وتصحيح الأخطاء والنشر على Google Play}

\chapterintro{صنعت التطبيق، وماذا بعد؟ يجيب هذا الفصل على ثلاث مسائل للانتقال من «يعمل» إلى «جاهز للتسليم»: كيف تصحح الأخطاء بكفاءة، وكيف تكتب اختبارات تمنع عودة المشاكل، وكيف ترفع تطبيقك إلى Google Play لينزّله العالم. والنشر ليس مجرد طقس — إنه الفاصل الذي يحوّل عملاً منجزاً إلى منتج حقيقي.}

\booksection{أولاً: أدوات التصحيح الثلاث — Logcat ونقاط التوقف وفاحص واجهة المستخدم}

\begin{center}
\begin{tabularx}{\linewidth}{|X|X|X|}
\hline
\textbf{الأداة} & \textbf{ماذا تفعل} & \textbf{كيف تستخدمها} \tabularnewline
\hline
Logcat & عرض سجلات النظام أو سجلاتك الخاصة & اكتب Log.d("TAG", "القيمة=\$x") في الكود ثم ابحث عن TAG في اللوحة للتصفية \tabularnewline
التصحيح بنقاط التوقف & إيقاف البرنامج عند سطر معيّن ومشاهدة قيم المتغيرات & انقر بجانب رقم السطر فتظهر نقطة حمراء، ثم شغّل بوضع Debug وتقدّم خطوة بخطوة \tabularnewline
Layout Inspector & فحص بنية الواجهة أثناء التشغيل & من القائمة Tools Layout Inspector لرؤية شجرة المكونات مباشرة \tabularnewline
\hline
\end{tabularx}
\end{center}

\begin{center}\small\color{gray}أدوات التصحيح الثلاث\end{center}

\begin{tcbinfo}{نصيحة: كيف تقرأ رسائل الخطأ؟}
بعد الانهيار يظهر في Logcat تتبع أحمر (Stacktrace)؛ ابحث من الأسفل إلى الأعلى عن أول سطر يبدأ باسم حزمتك (مثل com.example.yourapp)، فهذا السطر هو مكان الحادث غالباً. اقرأ أولاً «نوع الاستثناء» (مؤشر فارغ؟ تحويل نوع؟) ثم «at ... ملفك: رقم السطر». بهذه الحيلة تحل 80\% من الانهيارات.

\end{tcbinfo}

\booksection{ثانياً: الاختبارات — تأمين على الكود}

الاختبار اليدوي لا يُعتمد — تعدّل ميزة فتنكسر أخرى بصمت. الاختبارات هي «دواء الندم» و«السترة الواقية» معاً. تنقسم اختبارات Android إلى ثلاث طبقات، وعلى المبتدئ إتقان الطبقة الأولى أولاً:

\begin{center}
\begin{tabularx}{\linewidth}{|X|X|X|X|}
\hline
\textbf{الطبقة} & \textbf{ماذا تختبر} & \textbf{أين توضع} & \textbf{السرعة} \tabularnewline
\hline
اختبارات الوحدة & منطق خالص (حساب وتحقق) & app/src/test/ & من رتبة الثواني \tabularnewline
اختبارات الأجهزة (Instrumented) & سلوك الواجهة على جهاز حقيقي & app/src/androidTest/ & من رتبة الدقائق \tabularnewline
الاختبار اليدوي & إحساس وتجربة إنسان حقيقي & —— & لا غنى عنه أبداً \tabularnewline
\hline
\end{tabularx}
\end{center}

\begin{center}\small\color{gray}طبقات الاختبار الثلاث\end{center}

\begin{lstlisting}[language=Kotlin]
// استخرج منطق "تحويل نص المبلغ إلى رقم" كدالة نقية
fun parseAmount(input: String): Double? {
    val v = input.trim().toDoubleOrNull()
    return if (v != null && v > 0) v else null
}

// اختبرها (JUnit4)
@Test
fun emptyInputReturnsNull() {
    assertEquals(null, parseAmount(""))
}
@Test
fun negativeReturnsNull() {
    assertEquals(null, parseAmount("-5"))
}
@Test
fun normalValueParses() {
    assertEquals(12.5, parseAmount("12.5"))
}
\end{lstlisting}

التحقق من المبالغ والحسابات الإحصائية في دفتر المصاريف يجب أن تحميها اختبارات الوحدة. لتشغيلها: انقر بالزر الأيمن على مجلد test ثم Run. إذا اجتازت الاختبارات كلها بلون أخضر فعدّل كودك مطمئناً.

\booksection{ثالثاً: التوقيع والتجزئة — وضع «الختم» على التطبيق}

يشترط Android أن تكون كل حزمة تثبيت موقّعة — كالختم الرسمي الذي يثبت أن «هذه الحزمة مني ولم يُعبث بها». أثناء التصحيح يستخدم Android Studio توقيع debug تلقائياً، أما النشر فيتطلب توقيع release خاصاً بك. الخطوات:

\begin{enumerate}
  \item من القائمة Build Generate Signed Bundle / APK؛
  \item إنشاء ملف keystore (مخزن المفاتيح) — اضبط كلمة مرور واحفظه بعناية؛
  \item اختر صيغة release مع AAB ويكتمل البناء.
\end{enumerate}

\textbf{تحذير من فخ: ضياع keystore يعني ضياع التطبيق}\par

تحديث التطبيق بعد نشره يتطلب التوقيع نفسه، وإلا رفض المتجر التحديث. وإذا ضاع ملف keystore وكلمة المرور فلا سبيل إلا إلى تطبيق جديد يُنشر من الصفر — ولن يلتحق به المستخدمون القدامى. احرص على أخذ نسخة احتياطية من keystore وكلمة المرور في مكان آمن (سحابة مشفرة + دفتر ملاحظات)، فهذا من أغلى الدروس التي يدفعها مطورو Android.

\booksection{رابعاً: خطوات النشر على Google Play كاملة}

تسجيل حساب مطوّر

تجهيز مواد المتجر

رفع ملف AAB

ملء نموذج أمان البيانات

المراجعة (ساعات إلى أيام)

النشر

\begin{center}\small\color{gray}الشكل 6: ست خطوات للنشر على Google Play\end{center}

\begin{itemize}
  \item التسجيل: رسوم لمرة واحدة قدرها 25 دولاراً، مع حساب Google عامل؛
  \item المواد: أيقونة التطبيق (512×512)، ولقطات شاشة (اثنتان على الأقل)، ووصف الوظائف، ورابط سياسة الخصوصية (إلزامي عند جمع البيانات)؛
  \item استمارة أمان البيانات: صرّح بصدق عن البيانات التي تجمعها — دفتر المصاريف غير المتصل بالإنترنت يمكنه الإعلان عن «صفر جمع»؛
  \item المراجعة: النشر الأول للحسابات الجديدة يخضع غالباً لمراجعة بشرية تستغرق من ساعات إلى أيام، وإن رُفض فعدّل وفق الملاحظات.
\end{itemize}

ملاحظة للمستخدمين داخل الصين: يتطلب Google Play بيئة شبكة خاصة، ويتم توزيع التطبيقات في البر الرئيسي الصيني عادة عبر المتاجر المحلية (Huawei وXiaomi وTencent MyApp وغيرها). عملية النشر في المتاجر المحلية مشابهة، لكن لكل متجر ذوقه في المراجعة — ويمكن توزيع ملف AAB/APK نفسه على متاجر متعددة. وللتدريب، يكفيك النجاح مع متجر واحد.

\booksection{خامساً: إعداد مواد المتجر — الانطباع الأول هو تحميلاتك}

صفحة المتجر هي واجهة العرض التي يراها المستخدم أثناء تجوله. ثلاث مواد تحدد الانطباع الأول:

\begin{itemize}
  \item الأيقونة: بمقاس 512×512، ولا تلجم بلقطة شاشة. إن لم تكن مصمماً فولّد شكلاً هندسياً بسيطاً بأداة ذكاء اصطناعي، والمعيار أن تتعرف عليه بنظرة واحدة حتى بعد صغره على الشاشة الرئيسية للهاتف؛
  \item لقطات الشاشة: اثنتان على الأقل، ويفضّل من 4 إلى 6 تعرض أهم سيناريوهات الوظائف؛ مع جملة وصف بسيطة لكل لقطة (عبارة «سجّل مصروفاً في ثانية» أفضل بعشرة أضعاف من «حل متكامل ومحترف لإدارة الحسابات»)؛
  \item الوصف: أوضح في الجملة الأولى «ما هذا + لمن هو»، ولا تكدّس عبارات مثل «الأكثر احترافاً والأقوى» — فترتيب البحث في المتجر يعتمد أساساً على الكلمات المفتاحية ونسبة التحويل إلى تحميل، والصدق أكثر فاعلية من المبالغة.
\end{itemize}

\begin{tcbinfo}{نصيحة: قارن بنظرة إلى المنافسين}
قبل النشر تجول بين العشرة الأوائل من التطبيقات المشابهة في المتجر والتقط لقطات شاشة للمقارنة: ما نمط أيقوناتها؟ ما الميزات التي تعرضها اللقطات؟ كيف تكتب العناوين؟ المحاكاة ليست سرقة — تعلّم القوالب وضع محتواك أنت، فهذا أسرع درس في الإطلاق لأي منتج جديد.

\end{tcbinfo}

\booksection{خامساً: قائمة مختصرة لتحسين الأداء والإبلاغ عن الانهيارات}

قبل نشر التطبيق، تستحق هاتان المسألتان نصف ساعة لكل منهما:

\booksubsection{أربع نقاط لفحص الأداء}

\begin{enumerate}
  \item استخدم LazyColumn للقوائم: لا تعرض قائمة طويلة أبداً بـ Column العادي — فالتحميل الكسول هو الدرس الأول في أداء Android؛
  \item لا تدرج الصور بمقاسها الأصلي: استخدم مكتبة الصور Coil لتصغيرها وتخزينها مؤقتاً تلقائياً — صورة واحدة بحجم 5 ميغابايت كافية لتعطيل الأجهزة الضعيفة؛
  \item تجنّب الأعمال الثقيلة في الخيط الرئيسي: نفّذ قاعدة البيانات والشبكة والملفات الكبيرة داخل coroutines؛
  \item قِس باستخدام Profiler: محلّل الأداء المدمج في Android Studio يعرض الذاكرة والمعالج واستهلاك البطارية — قِس أولاً ثم حسّن، ولا تعتمد على الإحساس.
\end{enumerate}

\booksubsection{الإبلاغ عن الانهيارات}

عدم انهيار التطبيق أثناء اختبارك لا يعني أنه لن ينهار على هواتف المستخدمين. قبل النشر ادمج Firebase Crashlytics (رسمي من Google ومجاني): بضعة أسطر إعداد تكفي ليُرسل تتبع الانهيار تلقائياً عند حدوثه، فترى في لوحة التحكم «أي طراز، وأي إصدار، وأي سطر انهار». والنشر دون مراقبة للانهيارات كقيادة السيارة وعينان مغمضتان.

\begin{tcbpractice}{لنطبّق: سلّم دفتر المصاريف مرة واحدة}
(1) أضف 3 اختبارات وحدة لمنطق parseAmount وما شابهه؛ (2) اصطد بنفسك عبر Logcat انهياراً صنّعته ثم أصلحه؛ (3) أنشئ ملف APK موقّعاً بنسخة release وأرسله لصديق ليثبته (هذه هي متعة «العمل المنجز»)؛ (4) وإن سمحت طاقتك فسجّل حساب مطوّر وقدّم دفتر المصاريف فعلياً للنشر. الثلاث الأولى تنجزها الليلة.

\end{tcbpractice}

\begin{tcbsummary}{خلاصة الفصل}
\begin{itemize}
  \item أدوات التصحيح الثلاث: Logcat للسجلات، ونقاط التوقف للمتغيرات، وLayout Inspector للواجهة
  \item قراءة الانهيار: ابحث من أسفل التتبع إلى أعلاه عن أول سطر يحمل اسم حزمتك
  \item طبقات الاختبار الثلاث: الوحدة (من رتبة الثواني، تعلّمها أولاً) / الأجهزة / اليدوي
  \item النشر يتطلب توقيعاً: احفظ نسخة احتياطية من keystore وكلمة المرور، فضياعهما يعني «تطبيقاً بغير صاحبه»
  \item ست خطوات للنشر: التسجيل المواد رفع AAB استمارة أمان البيانات المراجعة النشر
\end{itemize}

\end{tcbsummary}
